Generative artificial intelligence (AI) offers a transformative pathway to bridge this clinical gap by synthesizing PET biomarkers from routinely acquired MRI~\cite{dayarathna2024deep,lee2024synthesizing}. Initial studies, primarily using generative adversarial networks (GANs)~\cite{pan2018synthesizing,pan2020spatially,apostolopoulos2022applications} and, more recently, diffusion models~\cite{xie2024synthesizing,chen2025multi,theodorou2025mri2pet}, have established the feasibility of producing visually realistic PET images. Subsequent advances have sought to improve fidelity by incorporating domain knowledge, for instance, through guidance from large vision-language models (VLMs)~\cite{wang2025unisyn,liu2025direct}.  However, visual realism alone does not establish subject-specific accuracy for cross-sectional diagnosis or MCI stage discrimination, particularly in a heterogeneous condition such as AD~\cite{ossenkoppele2022amyloid}. This challenge stems from the inherent nature of MRI-to-PET synthesis as an ill-posed inverse problem; neurodegeneration visible on MRI is typically a downstream effect of prior molecular pathologies captured by PET~\cite{jagust2018imaging}. Consequently, individuals with similar structural MRI presentations at early disease stages can exhibit vastly different pathological burdens, leading to ambiguous mapping. We therefore investigate patient-specific clinical and tracer information as complementary conditioning for MRI-to-PET synthesis.
